\section{Limitations}
Despite the favorable throughput characteristics demonstrated throughout this study, ARS remains subject to several architectural limitations and trade-offs that influence its applicability across different workload classes and hardware environments.

\subsection{Memory Overhead}
The current Aero implementation relies on auxiliary scatter buffers, spatial key caches, histogram structures, and temporary refinement storage.
As a result, the framework requires an additional $O(n)$ memory footprint beyond the input dataset.
While this design enables contiguous memory writes, contention-free partitioning, and improved cache utilization, it increases total DRAM consumption relative to in-place comparison sorters such as Quicksort and PDQsort.
Consequently, ARS may be less suitable for memory-constrained environments where auxiliary allocation costs outweigh throughput gains.

\subsection{Duplicate-Dominated Workloads}
Spatial classification performs best when the projected key space exhibits sufficient dispersion.
Under highly duplicate-dominated or heavily clustered distributions, a large fraction of elements may collapse into a small number of spatial partitions.
This increases local refinement costs and reduces the effectiveness of occupancy balancing.
Although Aero's quantile-adaptive mapping improves robustness under skewed distributions, duplicate-aware comparison sorters may still outperform the current ARS implementation in extreme duplicate-heavy scenarios.

\subsection{Cache-Capacity Dependence}
The efficiency of ARS is strongly influenced by cache residency.
The framework assumes that refinement partitions remain sufficiently small to benefit from the processor cache hierarchy.
As dataset scale increases, some partitions may exceed available cache capacity, exposing additional DRAM latency and reducing throughput efficiency.
This limitation motivated the development of Exp A, which explored recursive cache-aware refinement strategies designed to preserve cache locality for very large datasets.
However, such mechanisms are not currently incorporated into the primary Aero implementation.

\subsection{Quantile Estimation Accuracy}
Aero relies on a sampled approximation of the input distribution to construct its quantile lookup table.
Although this approach substantially improves occupancy stability for non-uniform workloads, its effectiveness depends on the representativeness of the sampled data.
Rapidly changing distributions, adversarial inputs, or highly localized data characteristics may reduce the accuracy of the approximation and lead to less balanced partition occupancy.
Future adaptive sampling strategies may further improve robustness under such conditions.

\subsection{Scope of Evaluation}
The experimental evaluation presented in this work focuses on shared-memory multi-core systems.
Distributed execution, NUMA-aware partitioning strategies, heterogeneous accelerators such as GPUs, and large-scale cluster deployments remain outside the scope of the current study.
Consequently, the results should be interpreted primarily as evidence of ARS behavior within modern shared-memory environments rather than a universal characterization across all computing platforms.

\subsection{Streaming Scope}
The Exp D architecture demonstrates that spatial classification can be extended to support high-throughput streaming ingestion through epoch-based micro-batching and concurrent spatial consolidation.
However, the current design targets append-oriented workloads and does not address persistence, fault tolerance, distributed stream coordination, or recovery semantics.
These concerns are essential for production-grade stream processing systems and represent important directions for future investigation.

Overall, the limitations identified in this section do not invalidate the central contributions of ARS, but rather highlight the architectural trade-offs associated with transforming sorting into a hardware-conscious spatial classification problem.
The framework exchanges additional memory usage and distribution sensitivity for improved throughput, reduced branch dependence, and scalable parallel execution, reflecting a deliberate design choice tailored toward modern multi-core data processing workloads.
